% Данные преподавателя
\newcommand{\Dep}{25}                               % Кафедра
\newcommand{\teacherPos}{доцент, канд. техн. наук}  % Должность
\newcommand{\teacherName}{Н.В. Степанов}              % Имя преподавателя

% Данные лабораторной работы
\newcommand{\labNum}{4} % Номер
\newcommand{\labTitle}{Моделирование оптимального приемника дискретных сигналов в канале с аддитивным белым гауссовским шумом} % Название
\newcommand{\courseName}{Общая теория связи}              % Название дисциплины

% Данные студента
\newcommand{\studentGroup}{2454}          % Номер группы
\newcommand{\studentName}{А.С. Калашников}       % Имя студента
\newcommand{\cityYear}{2026}                 % Год

% Параметры титульного листа перенесены из lab.docx (размеры Word — bp).
% Настройки локальны: основной текст сохраняет полуторный интервал.
\newgeometry{left=30mm,right=10mm,top=20mm,bottom=20mm}
\thispagestyle{empty}
\begingroup
\SingleSpacing
\fontsize{12bp}{13.8bp}\selectfont
\setlength{\parindent}{0pt}
\setlength{\parskip}{0pt}
\setlength{\tabcolsep}{0pt}
\renewcommand{\arraystretch}{1}

{\centering
\makebox[\textwidth][c]{МИНИСТЕРСТВО НАУКИ И ВЫСШЕГО ОБРАЗОВАНИЯ РОССИЙСКОЙ ФЕДЕРАЦИИ}\par
\makebox[\textwidth][r]{федеральное государственное автономное образовательное учреждение высшего образования}\par
{\fontsize{10bp}{11.5bp}\selectfont
«САНКТ-ПЕТЕРБУРГСКИЙ ГОСУДАРСТВЕННЫЙ УНИВЕРСИТЕТ\\
АЭРОКОСМИЧЕСКОГО ПРИБОРОСТРОЕНИЯ»\par}
\vspace{24bp}
КАФЕДРА № \Dep\par}

\vspace{60bp}
ОТЧЕТ\\
ЗАЩИЩЕН С ОЦЕНКОЙ\par
\vspace{6bp}
{\fontsize{12bp}{20.7bp}\selectfont ПРЕПОДАВАТЕЛЬ\par}
\vspace{6bp}
% Ширины полей и промежутков из таблицы образца.
\begin{tabular}{@{}p{162.15bp}p{14.2bp}p{141.05bp}p{13.85bp}p{150.7bp}@{}}%
\centering\teacherPos & & & & \centering\teacherName\tabularnewline
\cline{1-1}\cline{3-3}\cline{5-5}
\centering\fontsize{10bp}{9bp}\selectfont должность, уч. степень, звание & &
\centering\fontsize{10bp}{9bp}\selectfont подпись, дата & &
\centering\fontsize{10bp}{9bp}\selectfont инициалы, фамилия\tabularnewline
\end{tabular}\par

\vspace{64.1bp}
{\centering
{\fontsize{14bp}{16.1bp}\selectfont ОТЧЕТ О ЛАБОРАТОРНОЙ РАБОТЕ №\labNum\par}
\vspace{36bp}
{\fontsize{14bp}{16.1bp}\selectfont\MakeUppercase{\labTitle}\par}
\vspace{36bp}
\mbox{по дисциплине: \courseName}\par
\vspace{13.8bp}
}

\vspace{84bp}
{\fontsize{12bp}{20.7bp}\selectfont РАБОТУ ВЫПОЛНИЛ\par}
\vspace{6bp}
% Процент после преамбулы исключает пробел перед буквой С.
\begin{tabular}{@{}p{108.35bp}p{86.6bp}p{11.8bp}p{131.95bp}p{11.8bp}p{131.45bp}@{}}%
СТУДЕНТ гр. № & \centering\studentGroup & & & & \centering\studentName\tabularnewline
\cline{2-2}\cline{4-4}\cline{6-6}
& & & \centering\fontsize{10bp}{9bp}\selectfont подпись, дата & &
\centering\fontsize{10bp}{9bp}\selectfont инициалы, фамилия\tabularnewline
\end{tabular}\par

\vfill
{\centering Санкт-Петербург \cityYear\par}
\endgroup
\restoregeometry
